% !TEX root = ./newergap.tex

\section{Outline of the key ingredients}\label{subclaim-sec}

...

We begin with a weak version of the Dickson-Hardy-Littlewood prime tuples conjecture \cite{hardy}, which (following Pintz \cite{pintz-polignac}) we refer to as $\DHL[k,j]$.  Recall that for any $k_0 \in \mathbb{N}$, an \emph{admissible $k$-tuple} is a tuple ${\mathcal H} = (h_1,\ldots,h_{k})$ of $k$ increasing integers $h_1 < \ldots < h_{k}$ which avoids at least one residue class $a_p\ (p) := \{ a_p + np: n \in \Z \}$ for every $p$.  For instance, $(0,2,6)$ is an admissible $3$-tuple, but $(0,2,4)$ is not.

For any $k \geq j \geq 2$, we let $\DHL[k,j]$ denote the following claim:

\begin{claim}[Weak Dickson-Hardy-Littlewood conjecture,
  {$\DHL[k,j]$}] For any admissible $k$-tuple ${\mathcal
    H}=(h_1,\ldots,h_{k})$ there
  exist infinitely many translates $n + {\mathcal H} =
  (n+h_1,\ldots,n+h_{k})$ of ${\mathcal H}$ which contain at least $j$
  primes.
\end{claim}

For any $k$, let $H(k)$ denote the minimal diameter $h_k-h_1$ of an admissible $k$-tuple; thus for instance $H(3)=6$.  It is clear that for any natural numbers $m \geq 1$ and $k \geq m+1$, the claim $\DHL[k,m+1]$ implies that $H_m \leq H(k)$.  We will therefore deduce Theorem \ref{main} from a number of claims of the form $\DHL[k,j]$.  More precisely, we have

\begin{theorem}\label{main-dhl} Unconditionally, we have the following claims:
\begin{itemize}
\item[(i)] $\DHL[50,2]$.   
\item[(ii)] $\DHL[\num{35146},3]$.
\item[(iii)] $\DHL[\num{1628943},4]$. 
\item[(iv)] $\DHL[\num{74487363},5]$.  
\item[(v)] $\DHL[\num{3400000000},6]$.
\item[(vi)] $\DHL[k,m+1]$ whenever $m \geq 1$ and $k \geq C\exp( (4 - \frac{52}{283}) m )$ for some sufficiently large absolute constant $C$.
\end{itemize}
Assume the Elliott-Halberstam conjecture $\EH[\theta]$ for all $0 < \theta < 1$.  Then we have the following improvements:
\begin{itemize}
\item[(vii)] $\DHL[54,3]$.
\item[(viii)] $\DHL[\num{5511},4]$.
\item[(ix)] $\DHL[\num{41558},5]$.
\item[(x)] $\DHL[\num{309661},6]$.
\item[(xi)] $\DHL[k,m]$ whenever $m \geq 1$ and $k \geq C\exp( 2 m )$ for some sufficiently large absolute constant $C$.
\end{itemize}
Assume the generalized Elliott-Halberstam conjecture $\GEH[\theta]$ for all $0 < \theta < 1$.  Then
\begin{itemize}
\item[(xii)] $\DHL[3,2]$.
\end{itemize}
\end{theorem}

Theorem \ref{main} then follows from Theorem \ref{main-dhl} and the following bounds on $H(k)$:

\begin{theorem}[Bounds on $H(k)$]\label{hk-bound}\ 
\begin{itemize}
\item[(xii)] $H(3)=6$.
\item[(i)] $H(50) = 246$.
\item[(vii)] $H(54) = 270$.
\item[(viii)] $H(\num{5511}) \leq \num{52116}$.
\item[(ii)] $H(\num{35146}) \leq \num{395106}$.
\item[(ix)] $H(\num{41558}) \leq \num{474226}$.
\item[(x)] $H(\num{309661}) \leq \num{4137854}$.
\item[(iii)] $H(\num{1628943}) \leq \num{24462654}$.
\item[(iv)] $H(\num{74487363}) \leq \num{1497901734}$.
\item[(v)] $H(\num{3400000000}) \leq \num{82575303678}$.
\item[(vi), (xi)] In the asymptotic limit $k \to \infty$, one has $H(k) \leq k \log k + k \log\log k - k + o(k)$.
\end{itemize}
\end{theorem}

The last part of Theorem \ref{hk-bound} arises from considering the admissible $k$-tuple constructed by taking the first $k$ primes after $k$; see e.g. \cite[(3.2)]{polymath8a}.  For the non-asymptotic cases we use a variety of numerical methods, which we review in Appendix \ref{tuples-sec}. 

The proof of Theorem \ref{main-dhl} follows the Goldston-Pintz-Y{\i}ld{\i}r{\i}m strategy that was also used in all previous progress on this problem (e.g. \cite{gpy}, \cite{mp}, \cite{zhang}, \cite{polymath8a}, \cite{maynard-new}), namely that of constructing a sieve function adapted to an admissible $k$-tuple with good properties.  More precisely, we set
$$ w := \log \log \log x,$$
and set
$$ W := \prod_{p \leq w} p,$$
and observe the crude bound
\begin{equation}\label{W-bound}
 W \ll \log \log^{O(1)} x
\end{equation}
and observe the following simple ``pigeonhole principle'' criterion for $\DHL[k,m+1]$ (cf. \cite[Lemma 4.1]{polymath8a}, though the normalization here is slightly different):

\begin{lemma}[Criterion for $\DHL$]\label{crit} Let $k \geq 2$ and $m \geq 1$ be fixed integers, and define the normalisation constant
\begin{equation}\label{bnorm}
 B := \frac{\phi(W)}{W} \log x.
\end{equation}
Suppose that for each fixed admissible $k$-tuple $(h_1,\dots,h_k)$
  and each residue class $b\ (W)$ such that $b+h_i$ is coprime to $W$ for all $i=1,\dots,k$, one can find a non-negative weight
  function $\nu \colon \N \to \R^+$, fixed quantities $\alpha > 0$ and $\beta_1,\dots,\beta_k \geq
  0$, such that one has the
  asymptotic upper bound
\begin{equation}\label{s1}
 \sum_{\substack{x \leq n \leq 2x\\ n = b\ (W)}} \nu(n) \leq (\alpha+o(1)) B^{-k} \frac{x}{W},
 \end{equation}
the asymptotic lower bound
\begin{equation}\label{s2}
  \sum_{\substack{x \leq n \leq 2x\\ n = b\ (W)}} \nu(n) \theta(n+h_i) \geq (\beta_i-o(1)) B^{1-k} \frac{x}{\phi(W)}
\end{equation}
for all $h_i \in {\mathcal H}$, and the key inequality
\begin{equation}\label{key}
\frac{\beta_1 + \dots + \beta_k}{\alpha} > m.
\end{equation}
Then $\DHL[k,m+1]$ holds.  
\end{lemma}

\begin{proof}
Let $(h_1,\ldots,h_{k})$ be a fixed
  admissible $k$-tuple.  Since it is admissible, there is at least
  one residue class $b\ (W)$ such that $(b+h_i,W)=1$ for all $h_i
  \in {\mathcal H}$.  For an arithmetic function $\nu$ as in the
  lemma, we consider the quantity
$$
N:=\sum_{\substack{x \leq n \leq 2x\\ n = b\ (W)}}
 \nu(n) \left(\sum_{i=1}^{k} \theta(n+h_i) - m\log 3x\right).
$$
Combining \eqref{s1} and \eqref{s2}, we obtain the lower bound
$$ 
N\geq (\beta_1+\dots+\beta_k-o(1)) B^{1-k} \frac{x}{\phi(W)} - (m\alpha+o(1)) B^{-k} \frac{x}{W} \log 3x.
$$ 
From \eqref{bnorm} and the crucial condition \eqref{key}, it follows that
$N>0$ if $x$ is sufficiently large.
\par
On the other hand, the sum
$$
\sum_{i=1}^{k} \theta(n+h_i) - m\log 3x
$$ 
can only be positive if $n+h_i$ is prime for \emph{at least} $m+1$
indices $i=1, \ldots, k$.  We conclude that, for all sufficiently
large $x$, there exists some integer $n \in [x,2x]$ such that $n+h_i$
is prime for at least $m+1$ values of $i=1,\ldots,k$. 
\par
Since $(h_1,\dots,h_k)$ is an arbitrary admissible $k$-tuple,
$\DHL[k,m+1]$ follows.
\end{proof}

The objective is then to construct non-negative weights $\nu$ whose associated ratio $\frac{\beta_1+\dots+\beta_k}{\alpha}$ is as large as possible.   Our sieve majorants will be a variant of the multidimensional Selberg sieves used in \cite{maynard-new}.  As with all Selberg sieves, the $\nu$ are constructed as the square of certain (signed) divisor sums.  The divisor sums we will use will be finite linear combinations of products of ``one-dimensional'' divisor sums.  More precisely, for any fixed smooth compactly supported function $F: [0,+\infty) \to \R$, define the divisor sum $\lambda_F: \Z \to \R$ by the formula
\begin{equation}\label{lambdaf-def}
\lambda_F(n) := \sum_{d|n} \mu(d) F( \log_x n )
\end{equation}
where
\begin{equation}\label{logx-def}
\log_x n:= \frac{\log n}{\log x}.
\end{equation}
One should think of $\lambda_F$ as a smoothed out version of the indicator function to numbers $n$ which are ``almost prime'' in the sense that they have no prime factors less than $x^\eps$ for some small fixed $\eps>0$; see Proposition \ref{almostprime} for a more rigorous version of this heuristic.

The functions $\nu$ we will use will take the form
$$ \nu(n) = \left( \sum_{j=1}^J c_j \lambda_{F_{j,1}}(n+h_1) \dots \lambda_{F_{j,k}}(n+h_k) \right)^2 $$
for various fixed coefficients $c_1,\dots,c_J \in \R$ and fixed smooth compactly supported functions $F_{j,i}: [0,+\infty) \to \R$ with $j=1,\dots,J$ and $i=1,\dots,k$.  (One can of course absorb the constant $c_j$ into one of the $F_{j,i}$ if one wishes.)   Informally, $\nu$ is a smooth restriction to those $n$ for which $n+h_1,\dots,n+h_k$ are all almost prime.

Clearly, $\nu$ is a (positive-definite) linear combination of functions of the form
$$ n \mapsto \prod_{i=1}^k \lambda_{F_i}(n+h_i) \lambda_{G_i}(n+h_i)$$
for various smooth functions $F_1,\dots,F_k,G_1,\dots,G_k: [0,+\infty) \to \R$.  The sum appearing in \eqref{s1} can thus be decomposed into linear combinations of sums of the form
\begin{equation}\label{sfg-1}
 \sum_{\substack{x \leq n \leq 2x\\ n = b\ (W)}} \prod_{i=1}^k \lambda_{F_i}(n+h_i) \lambda_{G_i}(n+h_i).
\end{equation}
Also, since from \eqref{lambdaf-def} we clearly have
\begin{equation}\label{lambdan-prime}
\lambda_F(n) = F(0)
\end{equation}
when $n$ is prime and $F$ is supported on $[0,1]$, the sum appearing in \eqref{s2} can be similarly decomposed into linear combinations of sums of the form
\begin{equation}\label{sfg-2}
 \sum_{\substack{x \leq n \leq 2x\\ n = b\ (W)}} \theta(n+h_i) \prod_{1 \leq i' \leq k; i \neq i'} \lambda_{F_{i'}}(n+h_{i'}) \lambda_{G_{i'}}(n+h_{i'}).
\end{equation}

To estimate the latter sums, we use the following asymptotic, proven in Section \ref{sieving-sec}.  For each compactly supported $F: [0,+\infty) \to \R$, let 
\begin{equation}\label{S-def}
S(F) := \sup \{ x \geq 0: F(x) \neq 0\}
\end{equation}
 denote the upper range of the support of $F$ (with the convention that $S(0)=0$).

\begin{theorem}[Asymptotic for prime sums]\label{prime-asym}  Let $k \geq 2$ be fixed, let $(h_1,\dots,h_k)$ be a fixed admissible $k$-tuple, and let $b\ (W)$ be such that $b+h_i$ is coprime to $W$ for each $i=1,\dots,k$.  Let $1 \leq i_0 \leq k$ be fixed, and for each $1 \leq i \leq k$ distinct from $i_0$, let $F_{i}, G_{i}: [0,+\infty) \to \R$ be fixed smooth compactly supported functions.  Assume one of the following hypotheses:
\begin{itemize}
\item[(i)]  (Elliott-Halberstam) There exists $0 < \vartheta < 1$ such that $\EH[\vartheta]$ holds, and such that 
\begin{equation}\label{fg-upper}
\sum_{1 \leq i \leq k; i \neq i_0} ( S(F_{i}) + S(G_{i}) ) < \theta.
\end{equation}
\item[(ii)]  (Motohashi-Pintz-Zhang) There exists $0 \leq \varpi < 1/4$ and $\delta > 0$ such that $\MPZ[\varpi,\delta]$ holds, and such that
\begin{equation}\label{fg-upper-alt}
\sum_{1 \leq i \leq k; i \neq i_0} ( S(F_{i}) + S(G_{i}) ) < \frac{1}{2} + 2 \varpi
\end{equation}
and
\begin{equation}
\sup_{1 \leq i \leq k; i \neq i_0} \Bigl\{S(F_{i}), S(G_{i}) \Bigr\}< \delta.
\end{equation}
\end{itemize}
Then we have
\begin{equation}\label{theta-oo}
 \sum_{\substack{x \leq n \leq 2x\\ n = b\ (W)}} \theta(n+h_i) \prod_{1 \leq i \leq k; i \neq i_0} \lambda_{F_{i}}(n+h_{i}) \lambda_{G_{i}}(n+h_{i}) = 
(c+o(1)) B^{1-k} \frac{x}{\phi(W)}
\end{equation}
where
$$ c := \prod_{1 \leq i \leq k; i \neq i_0} \left(\int_0^1 F'(t_{i}) G'_{i'}(t_{i})\ dt_{i}\right).$$
Here of course $F'$ denotes the derivative of $F$.
\end{theorem}


To estimate the former sums, we use the following asymptotic, also proven in Section \ref{sieving-sec}.

\begin{theorem}[Asymptotic for non-prime sums]\label{nonprime-asym}  Let $k \geq 1$ be fixed, let $(h_1,\dots,h_k)$ be a fixed admissible $k$-tuple, and let $b\ (W)$ be such that $b+h_i$ is coprime to $W$ for each $i=1,\dots,k$.  For each fixed $1 \leq i \leq k$, let $F_{i}, G_{i}: [0,+\infty) \to \R$ be fixed smooth compactly supported functions.  Assume one of the following hypotheses:
\begin{itemize}
\item[(i)]  (Trivial case)  One has
\begin{equation}\label{easy-upper}
\sum_{i=1}^k ( S(F_i) + S(G_i) ) < 1.
\end{equation}
\item[(ii)]  (Generalized Elliott-Halberstam) There exists $0 < \vartheta < 1$ and $i_0 \in \{1,\dots,k\}$ such that $\GEH[\vartheta]$ holds, and
\begin{equation}\label{eh-upper}
\sum_{1 \leq i \leq k; i \neq i_0} ( S(F_{i}) + S(G_{i}) ) < \vartheta.
\end{equation}
\end{itemize}
Then we have
\begin{equation}\label{lflg}
 \sum_{\substack{x \leq n \leq 2x\\ n = b\ (W)}} \prod_{i=1}^k \lambda_{F_{i}}(n+h_{i}) \lambda_{G_{i}}(n+h_{i}) = 
(c+o(1)) B^{-k} \frac{x}{W},
\end{equation}
where
\begin{equation}\label{c-def}
 c := \prod_{i=1}^k \left(\int_0^1 F'_i(t_i) G'_i(t_i)\ dt_i\right).
\end{equation}
\end{theorem}

A key point in (ii) is that no upper bound on $S(F_{i_0})$ or $S(G_{i_0})$ is required (although, as we will see in Section \ref{geh-case}, the result is a little easier to prove when one has $S(F_{i_0})+S(G_{i_0}) < 1$).  This flexibility in the $F_{i_0}, G_{i_0}$ functions will be particularly crucial to obtain part (xii) of Theorem \ref{main-dhl} and Theorem \ref{main}.

\begin{remark}  Theorems \ref{prime-asym}, \ref{nonprime-asym} can be viewed as probabilistic assertions of the following form: if $n$ is chosen uniformly at random from the set $\{ x \leq n \leq 2x: n = b\ (W)\}$, then the random variables $\Lambda(n+h_i)$ and $\lambda_{F_j}(n+h_j) \lambda_{G_j}(n+h_j)$ for $i,j=1,\dots,k$ have mean $(1+o(1)) \frac{W}{\phi(W)}$ and $(\int_0^1 F'_j(t) G'_j(t)\ dt + o(1)) B^{-1}$ respectively, and furthermore these random variables enjoy a limited amount of independence, except for the fact (as can be seen from \eqref{lambdan-prime}) that $\Lambda(n+h_i)$ and $\lambda_{F_i}(n+h_i) \lambda_{G_i}(n+h_i)$ are highly correlated.  Note though that we do not have asymptotics for any sum which involves two or more factors of $\Lambda$, as such estimates are of a difficulty at least as great as that of the twin prime conjecture.
\end{remark}



Theorems \ref{prime-asym}, \ref{nonprime-asym} may be combined with Lemma \ref{crit} to reduce the task of establishing estimates of the form $\DHL[k,m+1]$ to that of establishing certain variational problems.  For instance, in Section \ref{may-sec} we reprove the following result of Maynard \cite[Proposition 4.2]{maynard-new}:

\begin{theorem}[Sieving on the standard simplex]\label{maynard-thm}  Let $k \geq 2$ and $m \geq 1$ be fixed integers. For any fixed compactly supported square-integrable function $F: [0,+\infty)^k \to \R$, define the functionals
\begin{equation}\label{i-def}
I(F) := \int_{[0,+\infty)^k} F(t_1,\dots,t_k)^2\ dt_1 \dots t_k
\end{equation}
and
\begin{equation}\label{ji-def}
 J_i(F) := \int_{[0,+\infty)^{k-1}} \left(\int_0^\infty F(t_1,\dots,t_k)\ dt_i\right)^2 dt_1 \dots dt_{i-1} dt_{i+1} \dots dt_k
\end{equation}
for $i=1,\dots,k$, and let $M_k$ be the supremum
\begin{equation}\label{mk4}
 M_k := \sup \frac{\sum_{i=1}^k J_i(F)}{I(F)}
\end{equation}
over all square integrable functions $F$ that are supported on the simplex
$$ {\mathcal R}_k := \{ (t_1,\dots,t_k) \in [0,+\infty)^k: t_1+\dots+t_k \leq 1 \}$$
and are not identically zero (up to almost everywhere equivalence, of course).  Suppose that there is $0 < \vartheta < 1$ such that $\EH[\vartheta]$ holds, and such that
$$ M_k > \frac{2m}{\vartheta}.$$
Then $\DHL[k,m+1]$ holds.
\end{theorem}

Parts (vii)-(xi) of Theorem \ref{main-dhl} (and hence Theorem \ref{main}) are then immediate from the following results, proven in Section \ref{asymptotics-sec}:

\begin{theorem}[Lower bounds on $M_k$]\label{mlower}\ 
\begin{itemize}
\item[(vii)] $M_{54} > 4.00223$.
\item[(viii)] $M_{5511} > 6$.
\item[(ix)] $M_{41558} > 8$.
\item[(x)] $M_{309661} > 10$.
\item[(xi)] One has $M_k \geq \log k - C$ for all sufficiently large $k$ and an absolute constant $C$.
\end{itemize}
\end{theorem}

For sake of comparison, in \cite[Proposition 4.3]{maynard-new} it was shown that $M_5 > 2$, $M_{105} > 4$, and $M_k \geq \log k - 2\log\log k - 2$ for all sufficiently large $k$.  As remarked in that paper, the sieves used on the bounded gap problem prior to the work in \cite{maynard-new} would essentially correspond, in this notation, to the choice of functions $F$ of the special form $F(t_1,\dots,t_k) := f(t_1+\dots+t_k)$, which severely limits the size of the ratio in
\eqref{mk4} (in particular, the analogue of $M_k$ in this special case cannot exceed $4$, as shown in \cite{sound}).

In the converse direction, in ??? we will also show the upper bound $M_k \leq \frac{k}{k-1} \log k$ for all $k \geq 2$, which shown in particular that the bounds in (vii) and (xi) of the above theorem cannot be significantly improved.  We also remark that Theorem \ref{mlower}(vii) also gives a weaker version $\DHL[54,2]$ of Theorem \ref{main-dhl}(i). 

We also have a variant of Theorem \ref{maynard-thm} which can accept inputs of the form $\MPZ[\varpi,\delta]$:

\begin{theorem}[Sieving on a truncated simplex]\label{maynard-trunc}  Let $k \geq 2$ and $m \geq 1$ be fixed integers. Let $0<\varpi<1/4$ and $0 <\delta < 1/2$ be such that $\MPZ[\varpi,\delta]$ holds.  For any $\alpha>0$, let $M_k^{[\alpha]}$ be the supremum \eqref{mk4}, where the supremum now ranges over all square-integrable $F$ supported in the \emph{truncated} simplex
\begin{equation}\label{ttk}
 \{ (t_1,\dots,t_k) \in [0,\alpha]^k: t_1+\dots+t_k \leq 1 \}
\end{equation}
and are not identically zero.  If
$$ M_k^{[\frac{\delta}{1/4+\varpi}]} > \frac{m}{1/4+\varpi},$$
then $\DHL[k,m+1]$ holds.
\end{theorem}

The following variant of Theorem \ref{mlower}, when combined with Theorem \ref{mpz-poly-improv}, then allows one to use Theorem \ref{maynard-trunc} to establish parts (ii)-(vi) of Theorem \ref{main-dhl} (and hence Theorem \ref{main}):

\begin{theorem}[Lower bounds on $M_k^{[\alpha]}$]\label{mlower-var}\ 
\begin{itemize}
\item[(ii)] There exist $\delta,\varpi>0$ with $1080 \varpi + 330 \delta < 13$ and $M_{\num{35146}}^{[\frac{\delta}{1/4+\varpi}]} > \frac{2}{1/4+\varpi}$.
\item[(iii)] There exist $\delta,\varpi>0$ with $1080 \varpi + 330 \delta < 13$ and $M_{\num{1628943}}^{[\frac{\delta}{1/4+\varpi}]} > \frac{3}{1/4+\varpi}$.
\item[(iv)] There exist $\delta,\varpi>0$ with $1080 \varpi + 330 \delta < 13$ and $M_{\num{74487363}}^{[\frac{\delta}{1/4+\varpi}]} > \frac{4}{1/4+\varpi}$.
\item[(v)] There exist $\delta,\varpi>0$ with $1080 \varpi + 330 \delta < 13$ and $M_{\num{3400000000}}^{[\frac{\delta}{1/4+\varpi}]} > \frac{5}{1/4+\varpi}$.
\item[(vi)] For all sufficiently large $k$, there exist $\delta,\varpi>0$ with $1080 \varpi + 330 \delta < 13$, $\varpi \geq \frac{13}{1080} - \frac{C}{\log k}$, and $M_k^{[\frac{\delta}{1/4+\varpi}]} \geq \log k - C$ for some absolute constant $C$.
\end{itemize}
\end{theorem}

Indeed, with Theorem \ref{mlower-var}(vi), Theorem \ref{mpz-poly-improv}, and Theorem\ref{maynard-trunc}, we see that $\DHL[k,m+1]$ holds whenever $k$ is sufficiently large and
$$ m \leq (\log k - C) (\frac{1}{4} + \frac{13}{1080} - \frac{C}{\log k})$$
which is in particular implied by
$$ m \leq \frac{\log k}{4 - \frac{52}{283}} - C'$$
for some absolute constant $C'$, giving Theorem \ref{main-dhl}(vi).

Now we give a more flexible variant of Theorem \ref{maynard-thm}, in which the support of $F$ is enlarged, at the cost of reducing the range of integration of the $J_i$.

\begin{theorem}[Sieving on an epsilon-enlarged simplex]\label{epsilon-trick}  Let $k \geq 2$ and $m \geq 1$ be fixed integers, and let $0 < \eps < 1$ be fixed also. For any fixed compactly supported square-integrable function $F: [0,+\infty)^k \to \R$, define the functionals
$$ J_{i,1-\eps}(F) := \int_{(1-\eps) \cdot {\mathcal R}_{k-1}} \left(\int_0^\infty F(t_1,\dots,t_k)\ dt_i\right)^2 dt_1 \dots dt_{i-1} dt_{i+1} \dots dt_k$$
for $i=1,\dots,k$, and let $M_{k,\eps}$ be the supremum
$$ M_{k,\eps} := \sup \frac{\sum_{i=1}^k J_{i,1-\eps}(F)}{I(F)}$$
over all square integrable functions $F$ that are supported on the simplex
$$ (1+\eps) \cdot {\mathcal R}_k = \{ (t_1,\dots,t_k) \in [0,+\infty)^k: t_1+\dots+t_k \leq 1+\eps \}$$
and are not identically zero.  Suppose that there is $0 < \vartheta < 1$, and assume one of the following two hypotheses:
\begin{itemize}
\item[(i)]  $\EH[\vartheta]$ holds, and $1+\eps < \frac{1}{\vartheta}$.
\item[(ii)] $\GEH[\vartheta]$ holds, and $\eps < \frac{1}{k-1}$.
\end{itemize}
If
$$ M_{k,\eps} > \frac{2m}{\vartheta}$$
then $\DHL[k,m+1]$ holds.
\end{theorem}

We prove this theorem in Section \ref{trick-sec}.  Part (i) of Theorem \ref{main-dhl}, and a weaker version $\DHL[4,2]$ of part (xii), then follow from Theorem \ref{bv-thm} and the following computations, proven in ???:

\begin{theorem}[Lower bounds on $M_{k,\eps}$]\label{mke-lower}\ 
\begin{itemize}
\item[(i)] $M_{50,1/25} > 4.0043$.
\item[(xii')] $M_{4,0.18} > 2.00235$.
\end{itemize}
\end{theorem}

We remark that computations in the proof of Theorem \ref{mke-lower}(xii') are simple enough that the bound may be checked by hand, without use of a computer.  The computations used to establish the full strength of Theorem \ref{main-dhl}(xii) are however significantly more complicated.

In fact, we may enlarge the support of $F$ further.  We give a version corresponding to part (ii) of Theorem \ref{epsilon-trick}; there is also a version corresponding to part (i) also, but we will not give it here as we will not have any use for it.

\begin{theorem}[Going beyond the epsilon enlargement]\label{epsilon-beyond}  Let $k \geq 2$ and $m \geq 1$ be fixed integers, let $0 <\vartheta < 1$ be such that $\GEH[\vartheta]$ holds, and let $0 < \eps < \frac{1}{k-1}$ be fixed also. Suppose that there is a fixed non-zero square-integrable function $F: [0,+\infty)^k \to \R$ supported in $\frac{k}{k-1} \cdot {\mathcal R}_k$, such that for $i=1,\dots,k$ one has the vanishing marginal condition
\begin{equation}\label{vanishing-marginal}
\int_0^\infty F(t_1,\dots,t_k)\ dt_i = 0
\end{equation}
 whenever $t_1,\dots,t_{i-1},t_{i+1},\dots,t_k \geq 0$ are such that
$$ t_1+\dots+t_{i-1}+t_{i+1}+\dots+t_k > 1+\eps.$$
Suppose that we also have the inequality
$$ \frac{\sum_{i=1}^k J_{i,\eps}(F)}{I(F)} > \frac{2m}{\vartheta}.$$
Then $\DHL[k,m+1]$ holds.
\end{theorem}

This theorem is proven in Section \ref{beyond-sec}.  By performing a delicate partitioning of a certain three-dimensional polytope and working with piecewise polynomial choices of $F$ on this polytope, it will give part (xii) of Theorem \ref{main-dhl}; see ???.

There are several other ways to combine Theorems \ref{prime-asym}, \ref{nonprime-asym} with equidistribution theorems on the primes to obtain results of the form $\DHL[k,m+1]$, but all of our attempts to do so either did not improve the numerology, or else were numerically infeasible to implement.